%% ============================================================
%%  Acronyms — article "EWC INT8 on MCU" (English version)
%%  Léonard Rivals — Sprint 40, CL-Embedded
%%
%%  Usage: \usepackage[acronym,toc]{glossaries} then \makeglossaries
%%         %% ============================================================
%%  Acronyms — article "EWC INT8 on MCU" (English version)
%%  Léonard Rivals — Sprint 40, CL-Embedded
%%
%%  Usage: \usepackage[acronym,toc]{glossaries} then \makeglossaries
%%         %% ============================================================
%%  Acronyms — article "EWC INT8 on MCU" (English version)
%%  Léonard Rivals — Sprint 40, CL-Embedded
%%
%%  Usage: \usepackage[acronym,toc]{glossaries} then \makeglossaries
%%         %% ============================================================
%%  Acronyms — article "EWC INT8 on MCU" (English version)
%%  Léonard Rivals — Sprint 40, CL-Embedded
%%
%%  Usage: \usepackage[acronym,toc]{glossaries} then \makeglossaries
%%         \input{glossary_en}
%%         \printglossary[type=\acronymtype]
%%
%%  Strict mirror of glossary_fr.tex: same 24 keys, same order.
%% ============================================================

%% Descriptions name other acronyms in plain words, never through \gls{...}:
%% an entry not cited in the text is not printed, and the link would then point
%% to a non-existent anchor.
\glsnoexpandfields

%% Style: "long form (SHORT)" at first use, description in the printed list.
\setacronymstyle{long-short-desc}

% ------------------------------------------------------------------
%  CONTINUAL LEARNING
% ------------------------------------------------------------------
\newacronym[description={Learning a sequence of tasks without access to past data, where
  the model must absorb the new without erasing the old}]%
  {CL}{CL}{Continual Learning}
\newacronym[description={Anti-forgetting method adding to the loss a quadratic penalty that
  restrains changes to the weights deemed important for past tasks, importance being
  estimated by the Fisher information. It stores no past data, which makes it compatible
  with a microcontroller}]%
  {EWC}{EWC}{Elastic Weight Consolidation}
\newacronym[description={Optimization algorithm updating the weights along the direction
  opposite to the gradient, estimated on a small batch of examples; this is the update run
  on-board in the online protocol}]%
  {SGD}{SGD}{Stochastic Gradient Descent}
\newacronym[description={Approach where a pretrained network is frozen in Flash and only a
  lightweight trainable layer is updated in RAM, sample by sample}]%
  {TinyOL}{TinyOL}{Tiny Online Learning}
\newacronym[description={Non-neural paradigm encoding signals into very high-dimensional
  vectors; learning reduces to elementary operations (addition, majority vote), without
  backpropagation}]%
  {HDC}{HDC}{Hyperdimensional Computing}
\newacronym[description={Machine learning on microcontrollers, characterized by three
  cumulative constraints: no external memory, no operating system, and less than 256\,kB of
  volatile memory}]%
  {TinyML}{TinyML}{Tiny Machine Learning}

% ------------------------------------------------------------------
%  QUANTIZATION
% ------------------------------------------------------------------
\newacronym[description={Quantization applied to an already trained model: simple to
  implement, but prone to losses when the scale is set without calibration or when tensor
  dynamics are wide. This is the regime whose collapse and recovery this article measures}]%
  {PTQ}{PTQ}{Post-Training Quantization}
\newacronym[description={Quantization simulated at every training step through
  \emph{fake-quant}, so that the model learns directly under reduced-precision constraints;
  costlier to train, it brings no decisive gain here over a properly calibrated PTQ}]%
  {QAT}{QAT}{Quantization-Aware Training}

% ------------------------------------------------------------------
%  HARDWARE TARGET AND MEMORY
% ------------------------------------------------------------------
\newacronym[description={Processor integrating core, memory and peripherals on a single
  chip, intended for self-contained embedded systems}]%
  {MCU}{MCU}{Microcontroller Unit}
\newacronym[description={Volatile memory holding the program's variables and buffers; the
  most critical resource of the target}]%
  {RAM}{RAM}{Random Access Memory}
\newacronym[description={Static volatile memory integrated in the microcontroller, fast but
  expensive in silicon area, hence available in very small amounts. It sets the memory
  budget of this work}]%
  {SRAM}{SRAM}{Static Random Access Memory}
\newacronym[description={Fast memory attached directly to the processor core, not reachable
  by some peripherals; on the target board it complements the main static memory}]%
  {CCM}{CCM}{Core Coupled Memory}
\newacronym[description={Circuit specialized in neural-network operations, able to run
  integer multiply-accumulates massively in parallel. Its absence on the target board
  explains why INT8 brings no speed gain there}]%
  {NPU}{NPU}{Neural Processing Unit}
\newacronym[description={Hardware floating-point unit; on the target board it makes FP32
  operations as fast as integer ones, which cancels the speed benefit expected from
  quantization}]%
  {FPU}{FPU}{Floating-Point Unit}
\newacronym[description={Instruction set applying one operation to several data items in a
  single cycle; this is the lever the integer path of this work lacks to turn a format gain
  into a time gain}]%
  {SIMD}{SIMD}{Single Instruction, Multiple Data}
\newacronym[description={Library of neural-network kernels optimized for Cortex-M cores,
  exploiting vector instructions; the direct route to lifting the latency paradox observed
  here}]%
  {CMSISNN}{CMSIS-NN}{Common Microcontroller Software Interface Standard --- Neural Network}
\newacronym[description={Instruction putting the core to sleep until the next interrupt;
  using it while waiting on the serial link lowers the idle current and makes the
  difference-based energy measurement usable}]%
  {WFI}{WFI}{Wait For Interrupt}

% ------------------------------------------------------------------
%  INSTRUMENTATION AND LINK
% ------------------------------------------------------------------
\newacronym[description={Debug unit of the Cortex-M core whose cycle counter times a code
  section to the cycle, then converts it to microseconds from the clock frequency. Every
  latency in this article comes from it}]%
  {DWT}{DWT}{Data Watchpoint and Trace}
\newacronym[description={Asynchronous serial link between board and host; it carries the
  samples frame by frame and reports the predictions}]%
  {UART}{UART}{Universal Asynchronous Receiver-Transmitter}
\newacronym[description={Checksum computed on every transmitted frame; a zero error rate
  guarantees that the compared predictions are those produced by the board and not the
  product of a corrupted transmission}]%
  {CRC}{CRC}{Cyclic Redundancy Check}

% ------------------------------------------------------------------
%  COMPUTE COST AND METRICS
% ------------------------------------------------------------------
\newacronym[description={Elementary operation multiplying two values and adding the product
  to an accumulator; the counting unit for the cost of a dense layer}]%
  {MAC}{MAC}{Multiply-Accumulate}
\newacronym[description={Elementary floating-point operation; one multiply-accumulate counts
  as two}]%
  {FLOP}{FLOP}{Floating-Point Operation}
\newacronym[description={Cost expressed at the bit level, proportional to the product of
  operand widths; it predicts a factor of 16 between FP32 and INT8 that the board does not
  deliver, illustrating the limits of an analytical metric}]%
  {BOP}{BOP}{Bit Operation}
\newacronym[description={Area under the ROC curve: the probability that an anomalous sample
  receives a higher score than a normal one. Independent of the decision threshold}]%
  {AUROC}{AUROC}{Area Under the ROC Curve}

%%         \printglossary[type=\acronymtype]
%%
%%  Strict mirror of glossary_fr.tex: same 24 keys, same order.
%% ============================================================

%% Descriptions name other acronyms in plain words, never through \gls{...}:
%% an entry not cited in the text is not printed, and the link would then point
%% to a non-existent anchor.
\glsnoexpandfields

%% Style: "long form (SHORT)" at first use, description in the printed list.
\setacronymstyle{long-short-desc}

% ------------------------------------------------------------------
%  CONTINUAL LEARNING
% ------------------------------------------------------------------
\newacronym[description={Learning a sequence of tasks without access to past data, where
  the model must absorb the new without erasing the old}]%
  {CL}{CL}{Continual Learning}
\newacronym[description={Anti-forgetting method adding to the loss a quadratic penalty that
  restrains changes to the weights deemed important for past tasks, importance being
  estimated by the Fisher information. It stores no past data, which makes it compatible
  with a microcontroller}]%
  {EWC}{EWC}{Elastic Weight Consolidation}
\newacronym[description={Optimization algorithm updating the weights along the direction
  opposite to the gradient, estimated on a small batch of examples; this is the update run
  on-board in the online protocol}]%
  {SGD}{SGD}{Stochastic Gradient Descent}
\newacronym[description={Approach where a pretrained network is frozen in Flash and only a
  lightweight trainable layer is updated in RAM, sample by sample}]%
  {TinyOL}{TinyOL}{Tiny Online Learning}
\newacronym[description={Non-neural paradigm encoding signals into very high-dimensional
  vectors; learning reduces to elementary operations (addition, majority vote), without
  backpropagation}]%
  {HDC}{HDC}{Hyperdimensional Computing}
\newacronym[description={Machine learning on microcontrollers, characterized by three
  cumulative constraints: no external memory, no operating system, and less than 256\,kB of
  volatile memory}]%
  {TinyML}{TinyML}{Tiny Machine Learning}

% ------------------------------------------------------------------
%  QUANTIZATION
% ------------------------------------------------------------------
\newacronym[description={Quantization applied to an already trained model: simple to
  implement, but prone to losses when the scale is set without calibration or when tensor
  dynamics are wide. This is the regime whose collapse and recovery this article measures}]%
  {PTQ}{PTQ}{Post-Training Quantization}
\newacronym[description={Quantization simulated at every training step through
  \emph{fake-quant}, so that the model learns directly under reduced-precision constraints;
  costlier to train, it brings no decisive gain here over a properly calibrated PTQ}]%
  {QAT}{QAT}{Quantization-Aware Training}

% ------------------------------------------------------------------
%  HARDWARE TARGET AND MEMORY
% ------------------------------------------------------------------
\newacronym[description={Processor integrating core, memory and peripherals on a single
  chip, intended for self-contained embedded systems}]%
  {MCU}{MCU}{Microcontroller Unit}
\newacronym[description={Volatile memory holding the program's variables and buffers; the
  most critical resource of the target}]%
  {RAM}{RAM}{Random Access Memory}
\newacronym[description={Static volatile memory integrated in the microcontroller, fast but
  expensive in silicon area, hence available in very small amounts. It sets the memory
  budget of this work}]%
  {SRAM}{SRAM}{Static Random Access Memory}
\newacronym[description={Fast memory attached directly to the processor core, not reachable
  by some peripherals; on the target board it complements the main static memory}]%
  {CCM}{CCM}{Core Coupled Memory}
\newacronym[description={Circuit specialized in neural-network operations, able to run
  integer multiply-accumulates massively in parallel. Its absence on the target board
  explains why INT8 brings no speed gain there}]%
  {NPU}{NPU}{Neural Processing Unit}
\newacronym[description={Hardware floating-point unit; on the target board it makes FP32
  operations as fast as integer ones, which cancels the speed benefit expected from
  quantization}]%
  {FPU}{FPU}{Floating-Point Unit}
\newacronym[description={Instruction set applying one operation to several data items in a
  single cycle; this is the lever the integer path of this work lacks to turn a format gain
  into a time gain}]%
  {SIMD}{SIMD}{Single Instruction, Multiple Data}
\newacronym[description={Library of neural-network kernels optimized for Cortex-M cores,
  exploiting vector instructions; the direct route to lifting the latency paradox observed
  here}]%
  {CMSISNN}{CMSIS-NN}{Common Microcontroller Software Interface Standard --- Neural Network}
\newacronym[description={Instruction putting the core to sleep until the next interrupt;
  using it while waiting on the serial link lowers the idle current and makes the
  difference-based energy measurement usable}]%
  {WFI}{WFI}{Wait For Interrupt}

% ------------------------------------------------------------------
%  INSTRUMENTATION AND LINK
% ------------------------------------------------------------------
\newacronym[description={Debug unit of the Cortex-M core whose cycle counter times a code
  section to the cycle, then converts it to microseconds from the clock frequency. Every
  latency in this article comes from it}]%
  {DWT}{DWT}{Data Watchpoint and Trace}
\newacronym[description={Asynchronous serial link between board and host; it carries the
  samples frame by frame and reports the predictions}]%
  {UART}{UART}{Universal Asynchronous Receiver-Transmitter}
\newacronym[description={Checksum computed on every transmitted frame; a zero error rate
  guarantees that the compared predictions are those produced by the board and not the
  product of a corrupted transmission}]%
  {CRC}{CRC}{Cyclic Redundancy Check}

% ------------------------------------------------------------------
%  COMPUTE COST AND METRICS
% ------------------------------------------------------------------
\newacronym[description={Elementary operation multiplying two values and adding the product
  to an accumulator; the counting unit for the cost of a dense layer}]%
  {MAC}{MAC}{Multiply-Accumulate}
\newacronym[description={Elementary floating-point operation; one multiply-accumulate counts
  as two}]%
  {FLOP}{FLOP}{Floating-Point Operation}
\newacronym[description={Cost expressed at the bit level, proportional to the product of
  operand widths; it predicts a factor of 16 between FP32 and INT8 that the board does not
  deliver, illustrating the limits of an analytical metric}]%
  {BOP}{BOP}{Bit Operation}
\newacronym[description={Area under the ROC curve: the probability that an anomalous sample
  receives a higher score than a normal one. Independent of the decision threshold}]%
  {AUROC}{AUROC}{Area Under the ROC Curve}

%%         \printglossary[type=\acronymtype]
%%
%%  Strict mirror of glossary_fr.tex: same 24 keys, same order.
%% ============================================================

%% Descriptions name other acronyms in plain words, never through \gls{...}:
%% an entry not cited in the text is not printed, and the link would then point
%% to a non-existent anchor.
\glsnoexpandfields

%% Style: "long form (SHORT)" at first use, description in the printed list.
\setacronymstyle{long-short-desc}

% ------------------------------------------------------------------
%  CONTINUAL LEARNING
% ------------------------------------------------------------------
\newacronym[description={Learning a sequence of tasks without access to past data, where
  the model must absorb the new without erasing the old}]%
  {CL}{CL}{Continual Learning}
\newacronym[description={Anti-forgetting method adding to the loss a quadratic penalty that
  restrains changes to the weights deemed important for past tasks, importance being
  estimated by the Fisher information. It stores no past data, which makes it compatible
  with a microcontroller}]%
  {EWC}{EWC}{Elastic Weight Consolidation}
\newacronym[description={Optimization algorithm updating the weights along the direction
  opposite to the gradient, estimated on a small batch of examples; this is the update run
  on-board in the online protocol}]%
  {SGD}{SGD}{Stochastic Gradient Descent}
\newacronym[description={Approach where a pretrained network is frozen in Flash and only a
  lightweight trainable layer is updated in RAM, sample by sample}]%
  {TinyOL}{TinyOL}{Tiny Online Learning}
\newacronym[description={Non-neural paradigm encoding signals into very high-dimensional
  vectors; learning reduces to elementary operations (addition, majority vote), without
  backpropagation}]%
  {HDC}{HDC}{Hyperdimensional Computing}
\newacronym[description={Machine learning on microcontrollers, characterized by three
  cumulative constraints: no external memory, no operating system, and less than 256\,kB of
  volatile memory}]%
  {TinyML}{TinyML}{Tiny Machine Learning}

% ------------------------------------------------------------------
%  QUANTIZATION
% ------------------------------------------------------------------
\newacronym[description={Quantization applied to an already trained model: simple to
  implement, but prone to losses when the scale is set without calibration or when tensor
  dynamics are wide. This is the regime whose collapse and recovery this article measures}]%
  {PTQ}{PTQ}{Post-Training Quantization}
\newacronym[description={Quantization simulated at every training step through
  \emph{fake-quant}, so that the model learns directly under reduced-precision constraints;
  costlier to train, it brings no decisive gain here over a properly calibrated PTQ}]%
  {QAT}{QAT}{Quantization-Aware Training}

% ------------------------------------------------------------------
%  HARDWARE TARGET AND MEMORY
% ------------------------------------------------------------------
\newacronym[description={Processor integrating core, memory and peripherals on a single
  chip, intended for self-contained embedded systems}]%
  {MCU}{MCU}{Microcontroller Unit}
\newacronym[description={Volatile memory holding the program's variables and buffers; the
  most critical resource of the target}]%
  {RAM}{RAM}{Random Access Memory}
\newacronym[description={Static volatile memory integrated in the microcontroller, fast but
  expensive in silicon area, hence available in very small amounts. It sets the memory
  budget of this work}]%
  {SRAM}{SRAM}{Static Random Access Memory}
\newacronym[description={Fast memory attached directly to the processor core, not reachable
  by some peripherals; on the target board it complements the main static memory}]%
  {CCM}{CCM}{Core Coupled Memory}
\newacronym[description={Circuit specialized in neural-network operations, able to run
  integer multiply-accumulates massively in parallel. Its absence on the target board
  explains why INT8 brings no speed gain there}]%
  {NPU}{NPU}{Neural Processing Unit}
\newacronym[description={Hardware floating-point unit; on the target board it makes FP32
  operations as fast as integer ones, which cancels the speed benefit expected from
  quantization}]%
  {FPU}{FPU}{Floating-Point Unit}
\newacronym[description={Instruction set applying one operation to several data items in a
  single cycle; this is the lever the integer path of this work lacks to turn a format gain
  into a time gain}]%
  {SIMD}{SIMD}{Single Instruction, Multiple Data}
\newacronym[description={Library of neural-network kernels optimized for Cortex-M cores,
  exploiting vector instructions; the direct route to lifting the latency paradox observed
  here}]%
  {CMSISNN}{CMSIS-NN}{Common Microcontroller Software Interface Standard --- Neural Network}
\newacronym[description={Instruction putting the core to sleep until the next interrupt;
  using it while waiting on the serial link lowers the idle current and makes the
  difference-based energy measurement usable}]%
  {WFI}{WFI}{Wait For Interrupt}

% ------------------------------------------------------------------
%  INSTRUMENTATION AND LINK
% ------------------------------------------------------------------
\newacronym[description={Debug unit of the Cortex-M core whose cycle counter times a code
  section to the cycle, then converts it to microseconds from the clock frequency. Every
  latency in this article comes from it}]%
  {DWT}{DWT}{Data Watchpoint and Trace}
\newacronym[description={Asynchronous serial link between board and host; it carries the
  samples frame by frame and reports the predictions}]%
  {UART}{UART}{Universal Asynchronous Receiver-Transmitter}
\newacronym[description={Checksum computed on every transmitted frame; a zero error rate
  guarantees that the compared predictions are those produced by the board and not the
  product of a corrupted transmission}]%
  {CRC}{CRC}{Cyclic Redundancy Check}

% ------------------------------------------------------------------
%  COMPUTE COST AND METRICS
% ------------------------------------------------------------------
\newacronym[description={Elementary operation multiplying two values and adding the product
  to an accumulator; the counting unit for the cost of a dense layer}]%
  {MAC}{MAC}{Multiply-Accumulate}
\newacronym[description={Elementary floating-point operation; one multiply-accumulate counts
  as two}]%
  {FLOP}{FLOP}{Floating-Point Operation}
\newacronym[description={Cost expressed at the bit level, proportional to the product of
  operand widths; it predicts a factor of 16 between FP32 and INT8 that the board does not
  deliver, illustrating the limits of an analytical metric}]%
  {BOP}{BOP}{Bit Operation}
\newacronym[description={Area under the ROC curve: the probability that an anomalous sample
  receives a higher score than a normal one. Independent of the decision threshold}]%
  {AUROC}{AUROC}{Area Under the ROC Curve}

%%         \printglossary[type=\acronymtype]
%%
%%  Strict mirror of glossary_fr.tex: same 24 keys, same order.
%% ============================================================

%% Descriptions name other acronyms in plain words, never through \gls{...}:
%% an entry not cited in the text is not printed, and the link would then point
%% to a non-existent anchor.
\glsnoexpandfields

%% Style: "long form (SHORT)" at first use, description in the printed list.
\setacronymstyle{long-short-desc}

% ------------------------------------------------------------------
%  CONTINUAL LEARNING
% ------------------------------------------------------------------
\newacronym[description={Learning a sequence of tasks without access to past data, where
  the model must absorb the new without erasing the old}]%
  {CL}{CL}{Continual Learning}
\newacronym[description={Anti-forgetting method adding to the loss a quadratic penalty that
  restrains changes to the weights deemed important for past tasks, importance being
  estimated by the Fisher information. It stores no past data, which makes it compatible
  with a microcontroller}]%
  {EWC}{EWC}{Elastic Weight Consolidation}
\newacronym[description={Optimization algorithm updating the weights along the direction
  opposite to the gradient, estimated on a small batch of examples; this is the update run
  on-board in the online protocol}]%
  {SGD}{SGD}{Stochastic Gradient Descent}
\newacronym[description={Approach where a pretrained network is frozen in Flash and only a
  lightweight trainable layer is updated in RAM, sample by sample}]%
  {TinyOL}{TinyOL}{Tiny Online Learning}
\newacronym[description={Non-neural paradigm encoding signals into very high-dimensional
  vectors; learning reduces to elementary operations (addition, majority vote), without
  backpropagation}]%
  {HDC}{HDC}{Hyperdimensional Computing}
\newacronym[description={Machine learning on microcontrollers, characterized by three
  cumulative constraints: no external memory, no operating system, and less than 256\,kB of
  volatile memory}]%
  {TinyML}{TinyML}{Tiny Machine Learning}

% ------------------------------------------------------------------
%  QUANTIZATION
% ------------------------------------------------------------------
\newacronym[description={Quantization applied to an already trained model: simple to
  implement, but prone to losses when the scale is set without calibration or when tensor
  dynamics are wide. This is the regime whose collapse and recovery this article measures}]%
  {PTQ}{PTQ}{Post-Training Quantization}
\newacronym[description={Quantization simulated at every training step through
  \emph{fake-quant}, so that the model learns directly under reduced-precision constraints;
  costlier to train, it brings no decisive gain here over a properly calibrated PTQ}]%
  {QAT}{QAT}{Quantization-Aware Training}

% ------------------------------------------------------------------
%  HARDWARE TARGET AND MEMORY
% ------------------------------------------------------------------
\newacronym[description={Processor integrating core, memory and peripherals on a single
  chip, intended for self-contained embedded systems}]%
  {MCU}{MCU}{Microcontroller Unit}
\newacronym[description={Volatile memory holding the program's variables and buffers; the
  most critical resource of the target}]%
  {RAM}{RAM}{Random Access Memory}
\newacronym[description={Static volatile memory integrated in the microcontroller, fast but
  expensive in silicon area, hence available in very small amounts. It sets the memory
  budget of this work}]%
  {SRAM}{SRAM}{Static Random Access Memory}
\newacronym[description={Fast memory attached directly to the processor core, not reachable
  by some peripherals; on the target board it complements the main static memory}]%
  {CCM}{CCM}{Core Coupled Memory}
\newacronym[description={Circuit specialized in neural-network operations, able to run
  integer multiply-accumulates massively in parallel. Its absence on the target board
  explains why INT8 brings no speed gain there}]%
  {NPU}{NPU}{Neural Processing Unit}
\newacronym[description={Hardware floating-point unit; on the target board it makes FP32
  operations as fast as integer ones, which cancels the speed benefit expected from
  quantization}]%
  {FPU}{FPU}{Floating-Point Unit}
\newacronym[description={Instruction set applying one operation to several data items in a
  single cycle; this is the lever the integer path of this work lacks to turn a format gain
  into a time gain}]%
  {SIMD}{SIMD}{Single Instruction, Multiple Data}
\newacronym[description={Library of neural-network kernels optimized for Cortex-M cores,
  exploiting vector instructions; the direct route to lifting the latency paradox observed
  here}]%
  {CMSISNN}{CMSIS-NN}{Common Microcontroller Software Interface Standard --- Neural Network}
\newacronym[description={Instruction putting the core to sleep until the next interrupt;
  using it while waiting on the serial link lowers the idle current and makes the
  difference-based energy measurement usable}]%
  {WFI}{WFI}{Wait For Interrupt}

% ------------------------------------------------------------------
%  INSTRUMENTATION AND LINK
% ------------------------------------------------------------------
\newacronym[description={Debug unit of the Cortex-M core whose cycle counter times a code
  section to the cycle, then converts it to microseconds from the clock frequency. Every
  latency in this article comes from it}]%
  {DWT}{DWT}{Data Watchpoint and Trace}
\newacronym[description={Asynchronous serial link between board and host; it carries the
  samples frame by frame and reports the predictions}]%
  {UART}{UART}{Universal Asynchronous Receiver-Transmitter}
\newacronym[description={Checksum computed on every transmitted frame; a zero error rate
  guarantees that the compared predictions are those produced by the board and not the
  product of a corrupted transmission}]%
  {CRC}{CRC}{Cyclic Redundancy Check}

% ------------------------------------------------------------------
%  COMPUTE COST AND METRICS
% ------------------------------------------------------------------
\newacronym[description={Elementary operation multiplying two values and adding the product
  to an accumulator; the counting unit for the cost of a dense layer}]%
  {MAC}{MAC}{Multiply-Accumulate}
\newacronym[description={Elementary floating-point operation; one multiply-accumulate counts
  as two}]%
  {FLOP}{FLOP}{Floating-Point Operation}
\newacronym[description={Cost expressed at the bit level, proportional to the product of
  operand widths; it predicts a factor of 16 between FP32 and INT8 that the board does not
  deliver, illustrating the limits of an analytical metric}]%
  {BOP}{BOP}{Bit Operation}
\newacronym[description={Area under the ROC curve: the probability that an anomalous sample
  receives a higher score than a normal one. Independent of the decision threshold}]%
  {AUROC}{AUROC}{Area Under the ROC Curve}
